\documentclass{amsart}
% For dealing with references we use the comment environment
\usepackage{verbatim}
\newenvironment{reference}{\comment}{\endcomment}
%\newenvironment{reference}{}{}
\newenvironment{slogan}{\comment}{\endcomment}
\newenvironment{history}{\comment}{\endcomment}

% For commutative diagrams we use Xy-pic
\usepackage[all]{xy}

% We use 2cell for 2-commutative diagrams.
\xyoption{2cell}
\UseAllTwocells

% We use multicol for the list of chapters between chapters
\usepackage{multicol}

% This is generally recommended for better output
\usepackage{lmodern}
\usepackage[T1]{fontenc}

% For cross-file-references

% Package for hypertext links:
\usepackage{hyperref}

% For any local file, say "hello.tex" you want to link to please

% Theorem environments.
%
\theoremstyle{plain}
\newtheorem{theorem}[subsection]{Theorem}
\newtheorem{proposition}[subsection]{Proposition}
\newtheorem{lemma}[subsection]{Lemma}

\theoremstyle{definition}
\newtheorem{definition}[subsection]{Definition}
\newtheorem{example}[subsection]{Example}
\newtheorem{exercise}[subsection]{Exercise}
\newtheorem{situation}[subsection]{Situation}

\theoremstyle{remark}
\newtheorem{remark}[subsection]{Remark}
\newtheorem{remarks}[subsection]{Remarks}

\numberwithin{equation}{subsection}

% Macros
%
\def\lim{\mathop{\mathrm{lim}}\nolimits}
\def\colim{\mathop{\mathrm{colim}}\nolimits}
\def\Spec{\mathop{\mathrm{Spec}}}
\def\Hom{\mathop{\mathrm{Hom}}\nolimits}
\def\Ext{\mathop{\mathrm{Ext}}\nolimits}
\def\SheafHom{\mathop{\mathcal{H}\!\mathit{om}}\nolimits}
\def\SheafExt{\mathop{\mathcal{E}\!\mathit{xt}}\nolimits}
\def\Sch{\mathit{Sch}}
\def\Mor{\mathop{\mathrm{Mor}}\nolimits}
\def\Ob{\mathop{\mathrm{Ob}}\nolimits}
\def\Sh{\mathop{\mathit{Sh}}\nolimits}
\def\NL{\mathop{N\!L}\nolimits}
\def\CH{\mathop{\mathrm{CH}}\nolimits}
\def\proetale{{pro\text{-}\acute{e}tale}}
\def\etale{{\acute{e}tale}}
\def\QCoh{\mathit{QCoh}}
\def\Ker{\mathop{\mathrm{Ker}}}
\def\Im{\mathop{\mathrm{Im}}}
\def\Coker{\mathop{\mathrm{Coker}}}
\def\Coim{\mathop{\mathrm{Coim}}}

% Boxtimes
%
\DeclareMathSymbol{\boxtimes}{\mathbin}{AMSa}{"02}

%
% Macros for moduli stacks/spaces
%
\def\QCohstack{\mathcal{QC}\!\mathit{oh}}
\def\Cohstack{\mathcal{C}\!\mathit{oh}}
\def\Spacesstack{\mathcal{S}\!\mathit{paces}}
\def\Quotfunctor{\mathrm{Quot}}
\def\Hilbfunctor{\mathrm{Hilb}}
\def\Curvesstack{\mathcal{C}\!\mathit{urves}}
\def\Polarizedstack{\mathcal{P}\!\mathit{olarized}}
\def\Complexesstack{\mathcal{C}\!\mathit{omplexes}}
% \Pic is the operator that assigns to X its picard group, usage \Pic(X)
% \Picardstack_{X/B} denotes the Picard stack of X over B
% \Picardfunctor_{X/B} denotes the Picard functor of X over B
\def\Pic{\mathop{\mathrm{Pic}}\nolimits}
\def\Picardstack{\mathcal{P}\!\mathit{ic}}
\def\Picardfunctor{\mathrm{Pic}}
\def\Deformationcategory{\mathcal{D}\!\mathit{ef}}

\setlength{\emergencystretch}{3em}
\begin{document}
\title[Stacks Project, Tag 07N1]{334 --- A p-power divided-power extension admits a de Rham homotopy inverse up to multiplication by p}
\maketitle
\noindent Copyright (C) 2005--2025 Johan de Jong. Stacks Project authors. Source: \href{https://stacks.math.columbia.edu/tag/07N1}{Tag 07N1}.

\noindent Editorial difficulty note (not author text). Difficulty H3: The proof constructs an explicit p-adically continuous integration operator on a divided-power basis for a pth-root extension. Its recursion must simultaneously control denominators, identify differential splitting and universal injectivity, and establish the factor and horizontality identities before the homotopy-inverse criterion applies..

\noindent Editorial provenance note (not author text). History: excerpt from revision \texttt{a0444\allowbreak 6e57e\allowbreak c1fbc\allowbreak 252a8\allowbreak 71afc\allowbreak ec775\allowbreak 2fb28\allowbreak 07b14}, crystalline.tex, lines 5177--5347. Reference presentation was adapted; original-author context and core support blocks were retained or added. The mathematical wording is unchanged. Licensed under GNU FDL 1.2 or later; no invariant sections or cover texts. See ../licenses/COPYING. Context blocks reproduce author definitions and supporting results; they are not additional counted problems.

\section{Crystals in quasi-coherent modules}
\label{section-quasi-coherent-crystals}

\noindent
In Situation \ref{situation-affine}.
Set $X = \Spec(C)$ and $S = \Spec(A)$. We are going to
classify crystals in quasi-coherent modules on $\text{Cris}(X/S)$.
Before we do so we fix some notation.

\medskip\noindent
Choose a polynomial ring $P = A[x_i]$ over $A$ and a surjection $P \to C$
of $A$-algebras with kernel $J = \Ker(P \to C)$. Set
\begin{equation}
\label{equation-D}
D = \lim_e D_{P, \gamma}(J) / p^eD_{P, \gamma}(J)
\end{equation}
for the $p$-adically completed divided power envelope.
This ring comes with a divided power ideal $\bar J$ and divided power
structure $\bar \gamma$, see Lemma \href{https://stacks.math.columbia.edu/tag/07KG}{lemma list properties (Tag 07KG)}.
Set $D_e = D/p^eD$ and denote $\bar J_e$ the image of $\bar J$ in $D_e$.
We will use the short hand
\begin{equation}
\label{equation-omega-D}
\Omega_D = \lim_e \Omega_{D_e/A, \bar\gamma} =
\lim_e \Omega_{D/A, \bar\gamma}/p^e\Omega_{D/A, \bar\gamma}
\end{equation}
for the $p$-adic completion of the module of divided power differentials,
see Lemma \href{https://stacks.math.columbia.edu/tag/07KK}{lemma differentials completion (Tag 07KK)}.
It is also the $p$-adic completion of
$\Omega_{D_{P, \gamma}(J)/A, \bar\gamma}$
which is free on $\text{d}x_i$, see
Lemma \href{https://stacks.math.columbia.edu/tag/07HW}{lemma module differentials divided power envelope (Tag 07HW)}.
Hence any element of $\Omega_D$ can be written uniquely as a sum
$\sum f_i\text{d}x_i$ with for all $e$ only finitely many $f_i$
not in $p^eD$. Moreover, the maps
$\text{d}_{D_e/A, \bar\gamma} : D_e \to \Omega_{D_e/A, \bar\gamma}$
fit together to define a divided power $A$-derivation
\begin{equation}
\label{equation-derivation-D}
\text{d} : D \longrightarrow \Omega_D
\end{equation}
on $p$-adic completions.

\medskip\noindent
We will also need the ``products $\Spec(D(n))$ of $\Spec(D)$'', see
Proposition \href{https://stacks.math.columbia.edu/tag/07JN}{proposition compute cohomology (Tag 07JN)} and its proof for an
explanation. Formally these are defined as follows. For $n \geq 0$ let
$J(n) = \Ker(P \otimes_A \ldots \otimes_A P \to C)$ where
the tensor product has $n + 1$ factors. We set
\begin{equation}
\label{equation-Dn}
D(n) = \lim_e
D_{P \otimes_A \ldots \otimes_A P, \gamma}(J(n))/
p^eD_{P \otimes_A \ldots \otimes_A P, \gamma}(J(n))
\end{equation}
equal to the $p$-adic completion of the divided power envelope.
We denote $\bar J(n)$ its divided power ideal and $\bar \gamma(n)$
its divided powers. We also introduce $D(n)_e = D(n)/p^eD(n)$ as well
as the $p$-adically completed module of differentials
\begin{equation}
\label{equation-omega-Dn}
\Omega_{D(n)} = \lim_e \Omega_{D(n)_e/A, \bar\gamma} =
\lim_e \Omega_{D(n)/A, \bar\gamma}/p^e\Omega_{D(n)/A, \bar\gamma}
\end{equation}
and derivation
\begin{equation}
\label{equation-derivation-Dn}
\text{d} : D(n) \longrightarrow \Omega_{D(n)}
\end{equation}
Of course we have $D = D(0)$. Note that the rings $D(0), D(1), D(2), \ldots$
form a cosimplicial object in the category of divided power rings.



\noindent
In the lemma below we will consider pairs $(M, \nabla)$ satisfying the
following conditions
\begin{enumerate}
\item
\label{item-complete}
$M$ is a $p$-adically complete $D$-module,
\item
\label{item-connection}
$\nabla : M \to M \otimes^\wedge_D \Omega_D$ is a connection, i.e.,
$\nabla(fm) = m \otimes \text{d}f + f\nabla(m)$,
\item
\label{item-integrable}
$\nabla$ is integrable
(see Remark \href{https://stacks.math.columbia.edu/tag/07I0}{remark connection (Tag 07I0)}), and
\item
\label{item-topologically-quasi-nilpotent}
$\nabla$ is {\it topologically quasi-nilpotent}: If we write
$\nabla(m) = \sum \theta_i(m)\text{d}x_i$ for some operators
$\theta_i : M \to M$, then for any $m \in M$ there are only finitely
many pairs $(i, k)$ such that $\theta_i^k(m) \not \in pM$.
\end{enumerate}
The operators $\theta_i$ are sometimes denoted
$\nabla_{\partial/\partial x_i}$ in the literature.
In the following lemma we construct a functor from crystals in quasi-coherent
modules on $\text{Cris}(X/S)$ to the category of such pairs. We will show
this functor is an equivalence in
Proposition \href{https://stacks.math.columbia.edu/tag/07JH}{proposition crystals on affine (Tag 07JH)}.



\section{Pulling back along purely inseparable maps}
\label{section-pull-back-along-pth-root}

\noindent
By an $\alpha_p$-cover we mean a morphism of the form
$$
X' = \Spec(C[z]/(z^p - c)) \longrightarrow \Spec(C) = X
$$
where $C$ is an $\mathbf{F}_p$-algebra and $c \in C$. Equivalently,
$X'$ is an $\alpha_p$-torsor over $X$. An {\it iterated
$\alpha_p$-cover}\footnote{This is nonstandard notation.}
is a morphism of schemes in characteristic
$p$ which is locally on the target a composition of finitely many
$\alpha_p$-covers. In this section we prove that pullback along
such a morphism induces a quasi-isomorphism on crystalline cohomology
after inverting the prime $p$. In fact, we prove a precise version
of this result. We begin with a preliminary lemma whose formulation
needs some notation.

\medskip\noindent
Assume we have a ring map $B \to B'$ and quotients $\Omega_B \to \Omega$ and
$\Omega_{B'} \to \Omega'$ satisfying the assumptions of
Remark \href{https://stacks.math.columbia.edu/tag/07KJ}{remark base change connection (Tag 07KJ)}.
Thus (\href{https://stacks.math.columbia.edu/tag/07PY}{equation base change map complexes (Tag 07PY)}) provides a
canonical map of complexes
$$
c_M^\bullet :
M \otimes_B \Omega^\bullet
\longrightarrow
M \otimes_B (\Omega')^\bullet
$$
for all $B$-modules $M$ endowed with integrable connection
$\nabla : M \to M \otimes_B \Omega_B$.

\medskip\noindent
Suppose we have $a \in B$, $z \in B'$, and a map $\theta : B' \to B'$
satisfying the following assumptions
\begin{enumerate}
\item
\label{item-d-a-zero}
$\text{d}(a) = 0$,
\item
\label{item-direct-sum}
$\Omega' = B' \otimes_B \Omega \oplus B'\text{d}z$; we write
$\text{d}(f) = \text{d}_1(f) + \partial_z(f) \text{d}z$
with $\text{d}_1(f) \in B' \otimes \Omega$ and $\partial_z(f) \in B'$
for all $f \in B'$,
\item
\label{item-theta-linear}
$\theta : B' \to B'$ is $B$-linear,
\item
\label{item-integrate}
$\partial_z \circ \theta = a$,
\item
\label{item-injective}
$B \to B'$ is universally injective (and hence $\Omega \to \Omega'$
is injective),
\item
\label{item-factor}
$af - \theta(\partial_z(f)) \in B$ for all $f \in B'$,
\item
\label{item-horizontal}
$(\theta \otimes 1)(\text{d}_1(f)) - \text{d}_1(\theta(f)) \in \Omega$
for all $f \in B'$ where
$\theta \otimes 1 : B' \otimes \Omega \to B' \otimes \Omega$
\end{enumerate}
These conditions are not logically independent.
For example, assumption (\ref{item-integrate}) implies
that $\partial_z(af - \theta(\partial_z(f))) = 0$.
Hence if the image of $B \to B'$ is the collection of
elements annihilated by $\partial_z$, then (\ref{item-factor})
follows. A similar argument can be made for condition (\ref{item-horizontal}).



\begin{situation}
\label{situation-affine}
Here $p$ is a prime number, $(A, I, \gamma)$ is a divided power
ring such that $A$ is a $\mathbf{Z}_{(p)}$-algebra, and $A \to C$ is a
ring map such that $IC = 0$ and such that $p$ is nilpotent in $C$.
\end{situation}

\begin{lemma}
\label{lemma-computation}
In Situation \ref{situation-affine}. Assume $D$ and $\Omega_D$ are as in
(\ref{equation-D}) and (\ref{equation-omega-D}).
Let $\lambda \in D$. Let $D'$ be the $p$-adic completion of
$$
D[z]\langle \xi \rangle/(\xi - (z^p - \lambda))
$$
and let $\Omega_{D'}$ be the $p$-adic completion of the module of
divided power differentials of $D'$ over $A$. For any pair $(M, \nabla)$
over $D$ satisfying (\ref{item-complete}), (\ref{item-connection}),
(\ref{item-integrable}), and (\ref{item-topologically-quasi-nilpotent})
the canonical map of complexes (\href{https://stacks.math.columbia.edu/tag/07PY}{equation base change map complexes (Tag 07PY)})
$$
c_M^\bullet : M \otimes_D^\wedge \Omega^\bullet_D
\longrightarrow
M \otimes_D^\wedge \Omega^\bullet_{D'}
$$
has the following property: There exists a map $e_M^\bullet$
in the opposite direction such that both $c_M^\bullet \circ e_M^\bullet$
and $e_M^\bullet \circ c_M^\bullet$ are homotopic to multiplication by $p$.
\end{lemma}

\begin{proof}
We will prove this using Lemma \href{https://stacks.math.columbia.edu/tag/07Q7}{lemma find homotopy (Tag 07Q7)} with $a = p$.
Thus we have to find $\theta : D' \to D'$ and prove
(\ref{item-d-a-zero}), (\ref{item-direct-sum}), (\ref{item-theta-linear}),
(\ref{item-integrate}), (\ref{item-injective}), (\ref{item-factor}),
(\ref{item-horizontal}). We first collect some information about the rings
$D$ and $D'$ and the modules $\Omega_D$ and $\Omega_{D'}$.

\medskip\noindent
Writing
$$
D[z]\langle \xi \rangle/(\xi - (z^p - \lambda))
=
D\langle \xi \rangle[z]/(z^p - \xi - \lambda)
$$
we see that $D'$ is the $p$-adic completion of the free $D$-module
$$
\bigoplus\nolimits_{i = 0, \ldots, p - 1}
\bigoplus\nolimits_{n \geq 0}
z^i \xi^{[n]} D
$$
where $\xi^{[0]} = 1$.
It follows that $D \to D'$ has a continuous $D$-linear section, in particular
$D \to D'$ is universally injective, i.e., (\ref{item-injective}) holds.
We think of $D'$ as a divided power algebra
over $A$ with divided power ideal $\overline{J}' = \overline{J}D' + (\xi)$.
Then $D'$ is also the $p$-adic completion of the divided power envelope
of the ideal generated by $z^p - \lambda$ in $D$, see
Lemma \href{https://stacks.math.columbia.edu/tag/07HC}{lemma describe divided power envelope (Tag 07HC)}. Hence
$$
\Omega_{D'} = \Omega_D \otimes_D^\wedge D' \oplus D'\text{d}z
$$
by Lemma \href{https://stacks.math.columbia.edu/tag/07HW}{lemma module differentials divided power envelope (Tag 07HW)}.
This proves (\ref{item-direct-sum}). Note that (\ref{item-d-a-zero})
is obvious.

\medskip\noindent
At this point we construct $\theta$. (We wrote a PARI/gp script theta.gp
verifying some of the formulas in this proof which can be found in the
scripts subdirectory of the Stacks project.) Before we do so we compute
the derivative of the elements $z^i \xi^{[n]}$. We have
$\text{d}z^i = i z^{i - 1} \text{d}z$. For $n \geq 1$ we have
$$
\text{d}\xi^{[n]} =
\xi^{[n - 1]} \text{d}\xi =
- \xi^{[n - 1]}\text{d}\lambda + p z^{p - 1} \xi^{[n - 1]}\text{d}z
$$
because $\xi = z^p - \lambda$. For $0 < i < p$ and $n \geq 1$ we have
\begin{align*}
\text{d}(z^i\xi^{[n]})
& =
iz^{i - 1}\xi^{[n]}\text{d}z + z^i\xi^{[n - 1]}\text{d}\xi \\
& =
iz^{i - 1}\xi^{[n]}\text{d}z + z^i\xi^{[n - 1]}\text{d}(z^p - \lambda) \\
& =
- z^i\xi^{[n - 1]}\text{d}\lambda +
(iz^{i - 1}\xi^{[n]} + pz^{i + p - 1}\xi^{[n - 1]})\text{d}z \\
& =
- z^i\xi^{[n - 1]}\text{d}\lambda +
(iz^{i - 1}\xi^{[n]} + pz^{i - 1}(\xi + \lambda)\xi^{[n - 1]})\text{d}z \\
& =
- z^i\xi^{[n - 1]}\text{d}\lambda +
((i + pn)z^{i - 1}\xi^{[n]} + p\lambda z^{i - 1}\xi^{[n - 1]})\text{d}z
\end{align*}
the last equality because $\xi \xi^{[n - 1]} = n\xi^{[n]}$.
Thus we see that
\begin{align*}
\partial_z(z^i) & = i z^{i - 1} \\
\partial_z(\xi^{[n]}) & = p z^{p - 1} \xi^{[n - 1]} \\
\partial_z(z^i\xi^{[n]}) & =
(i + pn) z^{i - 1} \xi^{[n]} + p \lambda z^{i - 1}\xi^{[n - 1]}
\end{align*}
Motivated by these formulas we define $\theta$ by the rules
$$
\begin{matrix}
\theta(z^j)
& = & p\frac{z^{j + 1}}{j + 1}
& j = 0, \ldots p - 1, \\
\theta(z^{p - 1}\xi^{[m]})
& = & \xi^{[m + 1]}
& m \geq 1, \\
\theta(z^j \xi^{[m]})
& = &
\frac{p z^{j + 1} \xi^{[m]} - \theta(p\lambda z^j \xi^{[m - 1]})}{(j + 1 + pm)}
& 0 \leq j < p - 1, m \geq 1
\end{matrix}
$$
where in the last line we use induction on $m$ to define our choice of
$\theta$. Working this out we get (for $0 \leq j < p - 1$ and $1 \leq m$)
$$
\theta(z^j \xi^{[m]}) =
\textstyle{\frac{p z^{j + 1} \xi^{[m]}}{(j + 1 + pm)} -
\frac{p^2 \lambda z^{j + 1} \xi^{[m - 1]}}{(j + 1 + pm)(j + 1 + p(m - 1))} +
\ldots +
\frac{(-1)^m p^{m + 1} \lambda^m z^{j + 1}}
{(j + 1 + pm) \ldots (j + 1)}}
$$
although we will not use this expression below. It is clear that $\theta$
extends uniquely to a $p$-adically continuous $D$-linear map on $D'$.
By construction we have (\ref{item-theta-linear}) and (\ref{item-integrate}).
It remains to prove (\ref{item-factor}) and (\ref{item-horizontal}).

\medskip\noindent
Proof of (\ref{item-factor}) and (\ref{item-horizontal}).
As $\theta$ is $D$-linear and continuous it suffices to prove that
$p - \theta \circ \partial_z$,
resp.\ $(\theta \otimes 1) \circ \text{d}_1 - \text{d}_1 \circ \theta$
gives an element of $D$, resp.\ $\Omega_D$ when evaluated on the
elements $z^i\xi^{[n]}$\footnote{This can be done by direct computation:
It turns out that $p - \theta \circ \partial_z$ evaluated on
$z^i\xi^{[n]}$ gives zero except for $1$ which is mapped to $p$ and
$\xi$ which is mapped to $-p\lambda$. It turns out that 
$(\theta \otimes 1) \circ \text{d}_1 - \text{d}_1 \circ \theta$
evaluated on $z^i\xi^{[n]}$ gives zero except for $z^{p - 1}\xi$
which is mapped to $-\lambda$.}.
Set $D_0 = \mathbf{Z}_{(p)}[\lambda]$ and
$D_0' = \mathbf{Z}_{(p)}[z, \lambda]\langle \xi \rangle/(\xi - z^p + \lambda)$.
Observe that each of the expressions above is an element of
$D_0'$ or $\Omega_{D_0'}$. Hence it suffices to prove the result
in the case of $D_0 \to D_0'$. Note that $D_0$ and $D_0'$
are torsion free rings and that $D_0 \otimes \mathbf{Q} = \mathbf{Q}[\lambda]$
and $D'_0 \otimes \mathbf{Q} = \mathbf{Q}[z, \lambda]$.
Hence $D_0 \subset D'_0$ is the subring of elements annihilated
by $\partial_z$ and (\ref{item-factor})
follows from (\ref{item-integrate}), see the discussion directly preceding
Lemma \href{https://stacks.math.columbia.edu/tag/07Q7}{lemma find homotopy (Tag 07Q7)}. Similarly, we have
$\text{d}_1(f) = \partial_\lambda(f)\text{d}\lambda$ hence
$$
\left((\theta \otimes 1) \circ \text{d}_1 - \text{d}_1 \circ \theta\right)(f)
=
\left(\theta(\partial_\lambda(f)) - \partial_\lambda(\theta(f))\right)
\text{d}\lambda
$$
Applying $\partial_z$ to the coefficient we obtain
\begin{align*}
\partial_z\left(
\theta(\partial_\lambda(f)) - \partial_\lambda(\theta(f))
\right)
& =
p \partial_\lambda(f) - \partial_z(\partial_\lambda(\theta(f))) \\
& =
p \partial_\lambda(f) - \partial_\lambda(\partial_z(\theta(f))) \\
& =
p \partial_\lambda(f) - \partial_\lambda(p f) = 0
\end{align*}
whence the coefficient does not depend on $z$ as desired.
This finishes the proof of the lemma.
\end{proof}
\end{document}
